% =====================================================================
% Chapter 3 --- Related Work and Research Gap                 [cross-cutting]
% Based on the author's systematic literature review (DEFESA SLR):
% PRISMA 2020 + Kitchenham 5-criterion QA + Wohlin snowballing.
% Centrepiece: the uniqueness matrix and the "why nobody did this" argument.
% =====================================================================

This chapter surveys the literature relevant to the perception pipeline of
Chapter~\ref{cha:pipeline} and establishes the gap that the dissertation fills.
Rather than an ad-hoc narrative, the review follows a systematic protocol, whose
methodology is described first; the body of work is then organised thematically,
and the chapter closes by making the gap explicit and arguing why it has remained
open.

\section{Review Methodology}

\subsection{Protocol and Research Question}

The review follows the PRISMA~2020 reporting guideline \cite{page2021prisma},
the quality-assessment practice of Kitchenham and Charters
\cite{kitchenham2007slr}, and the snowballing guidelines of Wohlin
\cite{wohlin2014snowball}. It is driven by a single research question, framed in
the PICOC form:

\begin{quote}
\emph{Which CNN-based techniques have been applied for situational-awareness
perception (sense) of autonomous aerial vehicles in urban environments, using
camera and point-cloud fusion for 3D object detection and tracking?}
\end{quote}

The \emph{population} is autonomous aerial vehicles in urban environments; the
\emph{intervention} is CNN-based pipelines combining 2D detection with
point-cloud processing; the \emph{comparison} baselines are monocular, stereo,
and LiDAR-only methods; the \emph{outcome} is 3D detection and tracking accuracy
with real-time feasibility; and the \emph{context} is urban sense-and-avoid
(perception only).

\subsection{Sources, Search, and Selection}

Automated searches of OpenAlex, Semantic Scholar, and academic web search
engines were run with three query families --- (i) UAV sense-and-avoid with deep
perception, (ii) frustum- and PointNet-based 3D detection, and (iii)
LiDAR--camera fusion for aerial 3D detection --- and complemented by backward and
forward snowballing (one-hop) from four seed papers: PointNet
\cite{qi2017pointnet}, Frustum PointNets \cite{qi2018frustum}, the UAV3D
benchmark \cite{ye2024uav3d}, and a recent UAV 3D-perception survey. Studies were
included if published in 2017--2026 (the year PointNet appeared), addressing
perception (detection, tracking, or situational awareness) for UAVs or for
ground-vehicle techniques transferable to UAVs, and containing a deep-learning
component; works addressing only avoidance, planning, or control, lacking a
learning component, or in unrelated domains were excluded. Each retained study
was scored on the five Kitchenham criteria of Table~\ref{tab:qa}, on a
$0$/$0.5$/$1$ scale; following the protocol, studies scoring below $2.5$ were
demoted to background text and excluded from the synthesis. The process was
managed with Rayyan, Zotero, and Parsifal.

\begin{table}[H]
  \centering
  \caption{Quality-assessment criteria (Kitchenham), each scored $0$/$0.5$/$1$.}
  \label{tab:qa}
  \begin{tabular}{@{}cl@{}}
    \toprule
    Criterion & Question \\
    \midrule
    Q1 & Are the problem and objectives clearly stated? \\
    Q2 & Is the methodology described in enough detail to reproduce? \\
    Q3 & Are there quantitative experiments on an established benchmark? \\
    Q4 & Does the work address a component of the proposed pipeline? \\
    Q5 & Is the contribution clearly differentiated from prior art? \\
    \bottomrule
  \end{tabular}
\end{table}

\subsection{Outcome of the Search}

The automated searches returned roughly $150$ records; five thematic
gap-filling searches added $50$; and snowballing added $44$, for $244$ records
in total. After de-duplication, $183$ unique records remained. Title-and-abstract
screening retained $174$ studies (with $7$ deferred to full-text assessment and
$2$ cross-snowball duplicates removed). Quality assessment passed $173$ studies
(mean score $4.05/5$, median $4.5$; $51\%$ from top-tier venues); a single survey
fell below the cut-off and was demoted to background. These $173$ studies form
the basis of the synthesis that follows.

\section{Two-Dimensional Detection for Aerial Imagery}

The pipeline's first stage is a 2D detector. The one-stage YOLO family
\cite{redmon2016yolo}, which regresses bounding boxes and class scores in a
single pass, dominates aerial detection because of its real-time throughput; a
substantial body of work adapts it to the small objects, oblique viewpoints, and
scale variation characteristic of aerial imagery. These detectors are mature and
are treated here as a reliable building block rather than a research
contribution of this work.

\section{Frustum-Based 3D Detection in Autonomous Driving}

The core of the proposed pipeline descends from Frustum PointNets
\cite{qi2018frustum}, which lifts each 2D detection into a viewing frustum, crops
the corresponding point cloud, and applies a PointNet \cite{qi2017pointnet,
qi2017pointnetpp} to estimate an amodal 3D box. The idea spawned a rich lineage
in the driving domain: Frustum ConvNet \cite{wang2019fconvnet} replaces the
point-wise network with sliding-frustum convolutions, and Frustum-PointPillars
\cite{paigwar2021fpp} substitutes a pillar encoding for the PointNet backbone.
All of these methods are benchmarked on ground-vehicle datasets and have not been
applied to aerial platforms.

\section{Camera--LiDAR Fusion and Multi-View 3D Detection}

Beyond the frustum family, a broad literature fuses camera and LiDAR for 3D
detection. Region-level fusion methods such as AVOD \cite{ku2018avod} and
PointFusion \cite{xu2018pointfusion} combine image and point features without an
explicit frustum stage, while transformer- and bird's-eye-view-based detectors
such as DETR3D \cite{wang2022detr3d} and BEVFusion \cite{liu2023bevfusion}
operate on multi-camera or unified BEV representations. These approaches achieve
state-of-the-art accuracy on driving benchmarks but are LiDAR- or
multi-camera-centric and, again, are not deployed on UAVs.

\section{3D Perception and Sense-and-Avoid for UAVs}

On the aerial side, 3D perception is comparatively recent and follows different
design choices. The UAV3D benchmark \cite{ye2024uav3d} provides the first
large-scale aerial 3D dataset, but is a benchmark rather than a detection method
and is synthetic. Method papers favour LiDAR-centric or RGB-D pipelines:
LiDAR--camera obstacle detection for UAVs \cite{ma2021uavfusion}, aerial
LiDAR-based detection with unscented-Kalman-filter tracking
\cite{cherif2023aerial}, and airborne LiDAR sense-and-detect \cite{dolph2023airborne}.
Crucially, \emph{none} of these adopts the frustum-to-PointNet fusion that is
standard in driving --- they use voxel networks, RGB-D, or LiDAR alone.

\section{Multi-Object Tracking on Moving Platforms}

For the tracking stage, the literature divides into learning-based
appearance trackers and recursive-filter trackers. Filter-based multi-object
tracking on moving platforms --- the regime of this dissertation --- typically
couples a detector with a Kalman or extended Kalman filter; recent UAV trackers
such as DMTrack \cite{fu2025dmtrack} address association and motion but are
tracker-only, without an integrated 3D detection pipeline. The author's own prior
work \cite{delima2025fusion}, extended in Chapter~\ref{cha:tracking}, contributes
the recursive-filter tracking stage used here.

\section{Datasets for Aerial 3D Perception}

The scarcity of aerial 3D data is itself a structural obstacle. The driving
domain is served by KITTI \cite{geiger2012kitti}, nuScenes
\cite{caesar2020nuscenes}, and the Waymo Open Dataset \cite{sun2020waymo}, with
millions of annotated 3D boxes. For UAVs, large-scale 2D benchmarks such as
VisDrone \cite{zhu2021visdrone} exist, but the only large-scale 3D aerial
benchmark, UAV3D \cite{ye2024uav3d}, is one year old and synthetic. This
motivates the synthetic dataset built with AirSim \cite{shah2018airsim} in
Chapter~\ref{cha:airsim}.

\section{The Research Gap}

\subsection{Uniqueness Matrix}

Table~\ref{tab:uniqueness} maps the closest adjacent works against the seven
attributes that together define this dissertation's target: a drone target,
a frustum stage, a PointNet 3D-box estimator, camera--LiDAR fusion, an integrated
tracker, an urban environment, and a sense-and-avoid application. The works
divide into two blocks. The first --- the frustum and fusion methods from driving
--- possesses the full detection machinery but never targets UAVs. The second ---
the UAV-specific methods --- targets drones but lacks the frustum-to-PointNet
pipeline. Of the nineteen closest works examined, \emph{none} satisfies all seven
attributes; only this dissertation fills every cell.

\begin{table}[H]
  \centering
  \footnotesize
  \caption{Uniqueness matrix. \checkmark{}: present and central; $\sim$: partial;
  $\times$: absent. Only the last row fills all seven attributes.}
  \label{tab:uniqueness}
  \begin{tabular}{@{}lccccccc@{}}
    \toprule
    Work & Drone & Frustum & PointNet & Fusion & Track & Urban & S\&A \\
    \midrule
    Frustum PointNets \cite{qi2018frustum} & $\times$ & \checkmark & \checkmark & \checkmark & $\times$ & $\sim$ & $\times$ \\
    Frustum ConvNet \cite{wang2019fconvnet} & $\times$ & \checkmark & \checkmark & \checkmark & $\times$ & $\sim$ & $\times$ \\
    Frustum-PointPillars \cite{paigwar2021fpp} & $\times$ & \checkmark & $\times$ & \checkmark & $\times$ & $\sim$ & $\times$ \\
    Multi-View Frustum & $\times$ & \checkmark & \checkmark & \checkmark & $\times$ & $\sim$ & $\times$ \\
    Frustum-Fusion (embedded) & $\times$ & \checkmark & \checkmark & \checkmark & $\times$ & $\sim$ & $\times$ \\
    PointFusion \cite{xu2018pointfusion} & $\times$ & $\times$ & \checkmark & \checkmark & $\times$ & $\sim$ & $\times$ \\
    AVOD \cite{ku2018avod} & $\times$ & $\times$ & $\times$ & \checkmark & $\times$ & $\sim$ & $\times$ \\
    DETR3D \cite{wang2022detr3d} & $\times$ & $\times$ & $\times$ & $\sim$ & $\times$ & $\sim$ & $\times$ \\
    BEVFusion \cite{liu2023bevfusion} & $\times$ & $\times$ & $\times$ & \checkmark & $\times$ & $\sim$ & $\times$ \\
    UAV3D benchmark \cite{ye2024uav3d} & \checkmark & $\times$ & $\times$ & $\sim$ & $\sim$ & \checkmark & $\sim$ \\
    UAV LiDAR--cam \cite{ma2021uavfusion} & \checkmark & $\times$ & $\times$ & \checkmark & $\times$ & \checkmark & \checkmark \\
    Aerial LiDAR + UKF \cite{cherif2023aerial} & \checkmark & $\times$ & $\times$ & $\times$ & \checkmark & \checkmark & $\times$ \\
    Airborne LiDAR \cite{dolph2023airborne} & \checkmark & $\times$ & $\times$ & $\times$ & $\times$ & $\sim$ & \checkmark \\
    DMTrack \cite{fu2025dmtrack} & \checkmark & $\times$ & $\times$ & $\times$ & \checkmark & $\sim$ & $\times$ \\
    \midrule
    \textbf{This dissertation} & \checkmark & \checkmark & \checkmark & \checkmark & \checkmark & \checkmark & \checkmark \\
    \bottomrule
  \end{tabular}
\end{table}

The gap can be stated in one sentence: \emph{no published work has transferred
the frustum-to-PointNet detection pipeline from autonomous driving to UAV
sense-and-avoid in urban environments}. The closest mention in the surveyed
literature is a single aspirational sentence in a frustum-sampling paper noting
that such methods ``show promise'' for drones --- with no implementation or
evaluation.

\subsection{Why This Has Not Been Done Yet}

A natural question is why, if the transfer is so direct, it has not been done.
The review identifies six independent reasons:

\begin{enumerate}
  \item \textbf{LiDAR cost and weight.} Car-grade LiDAR ($\approx 13$~kg,
  $\sim$\$75k) is prohibitive on small UAVs; light, affordable UAV-grade LiDAR
  only matured around 2020--2022, so the sensing side of the problem did not
  exist before then.
  \item \textbf{Aerial annotation scarcity.} There is no KITTI/nuScenes/Waymo for
  UAVs; the first large-scale aerial 3D benchmark, UAV3D
  \cite{ye2024uav3d}, is one year old and synthetic, making it hard to train
  point-cloud networks on aerial data.
  \item \textbf{Non-trivial domain shift.} Scale, viewpoint, motion model (free
  3D rotation and translation versus a car's 2D motion), sensor noise, and
  obstacle geometry all differ; naively retraining a driving detector on aerial
  data does not work.
  \item \textbf{Onboard compute constraints.} UAV-grade compute (Jetson-class) is
  an order of magnitude weaker than a car's; real-time frustum-plus-PointNet
  inference at that budget is itself an active research subarea.
  \item \textbf{Funding-priority bias.} Most autonomy funding has targeted ground
  vehicles; Urban Air Mobility became a serious commercial agenda only in the
  late 2010s.
  \item \textbf{Deferred tracker integration.} 3D-detection papers typically stop
  at single-frame detection and mark tracking as future work, whereas for
  sense-and-avoid the tracker is the safety-critical stage.
\end{enumerate}

Combinations of these reasons explain the timing: reasons~(1)--(2) mean the
infrastructure did not exist before $\sim$2022; reasons~(3)--(4) mean the
engineering is hard even with it; and reasons~(5)--(6) mean the incentives were
thin until Urban Air Mobility matured.

\section{Positioning of This Dissertation}

Against this landscape, the dissertation occupies the empty cell of
Table~\ref{tab:uniqueness}: it adapts the frustum-to-PointNet pipeline from
ground vehicles to urban UAV perception, integrates 2D detection, frustum
proposal, 3D box estimation, and recursive-filter tracking into one
situational-awareness pipeline, and addresses the data and compute obstacles
through a synthetic dataset and an efficient, rigorously analysed tracker. The
tracking stage --- the safety-critical component that the literature most often
defers --- is the most mature contribution and is developed first, in
Chapter~\ref{cha:tracking}.
